\textbf{PHẦN I. Câu trắc nghiệm nhiều phương án lựa chọn}\setcounter{ex}{0}}
\begin{ex}
Nghiệm của phương trình $\sin x = 1$ là
\choice
{$x = k\pi \ (k \in \mathbb{Z})$}
{$x = \dfrac{\pi}{2} + k\pi \ (k \in \mathbb{Z})$}
{$x = \dfrac{\pi}{2} + k2\pi \ (k \in \mathbb{Z})$}
{$x = -\dfrac{\pi}{2} + k2\pi \ (k \in \mathbb{Z})$}
\loigiai{Vì $\sin x = 1 \Rightarrow x = \dfrac{\pi}{2} + k2\pi, \ k \in \mathbb{Z}$. Đáp án đúng là C}
\end{ex}

\begin{ex}
Công thức nào sau đây sai?
\choice
{$\sin(a+b) = \sin a \cos b + \cos a \sin b$}
{$\cos(a-b) = \cos a \cos b + \sin a \sin b$}
{$\cos(a+b) = \cos a \cos b + \sin a \sin b$}
{$\sin(a-b) = \sin a \cos b - \cos a \sin b$}
\loigiai{Công thức sai là $\cos(a+b) = \cos a \cos b + \sin a \sin b$. Đúng phải là $\cos(a+b) = \cos a \cos b - \sin a \sin b$. Đáp án C}
\end{ex}

\begin{ex}
Cho $0 < \alpha < \dfrac{\pi}{2}$. Khẳng định nào sau đây đúng?
\choice
{$\cos\alpha < 0$}
{$\cot\alpha < 0$}
{$\sin\alpha > 0$}
{$\tan\alpha < 0$}
\loigiai{Vì $0 < \alpha < \dfrac{\pi}{2}$ nên $\sin\alpha > 0$. Đáp án đúng là C}
\end{ex} 

\begin{ex}
Góc có số đo $110^\circ$ đổi ra đơn vị radian là
\choice
{$110\pi$}
{$\dfrac{110\pi}{3}$}
{$\dfrac{11\pi}{18}$}
{$\dfrac{\pi}{8}$}
\loigiai{$110^\circ = 110 \times \dfrac{\pi}{180} = \dfrac{11\pi}{18}$. Đáp án C}
\end{ex}
\begin{ex}
Cho các dãy số sau. Dãy số nào là dãy số tăng?

\choice
{$1; 3; 5; 7; 9; \ldots$}
{$1; -\dfrac{1}{2}; \dfrac{1}{4}; -\dfrac{1}{8}; \dfrac{1}{16}; \ldots$}
{$2; 2; 2; 2; 2; \ldots$}
{$1; \dfrac{1}{2}; \dfrac{1}{4}; \dfrac{1}{8}; \dfrac{1}{16}; \ldots$}
\end{ex}

\begin{ex}
Cho cấp số nhân $(u_n)$ có số hạng đầu $u_1 = 3$ và công bội $q = 2$. Giá trị $u_4$ là
\choice
{$24$}
{$9$}
{$-32$}
{$6$}
\loigiai{$u_4 = u_1 \cdot q^{4-1} = 3 \cdot 2^3 = 24$. Đáp án A}
\end{ex}
\begin{ex}
Cho cấp số cộng $(u_n)$ thỏa mãn $u_1 = 4, \ u_2 = 10$.  
Công sai của cấp số cộng bằng
\choice
{$4$}
{$\dfrac{5}{2}$}
{$10$}
{$6$}
\loigiai{
	Công sai $d = u_2 - u_1 = 10 - 4 = 6$.  
	\textbf{Đáp án đúng: D}
}
\end{ex}

\begin{ex}
Cho dãy số $(u_n)$ biết $u_n = 2n + 1$.  
Số hạng $u_3$ bằng
\choice
{$7$}
{$3$}
{$8$}
{$6$}
\loigiai{
	$u_3 = 2 \times 3 + 1 = 7$.  
	\textbf{Đáp án đúng: A}
}
\end{ex}
\begin{ex}
Phương trình lượng giác $\tan x = \dfrac{\sqrt{3}}{3}$ có nghiệm là
\choice
{$x = -\dfrac{\pi}{6} + k\pi \ (k \in \mathbb{Z})$}
{$x = -\dfrac{\pi}{3} + k2\pi \ (k \in \mathbb{Z})$}
{$x = -\dfrac{\pi}{3} + k\pi \ (k \in \mathbb{Z})$}
{$x = \dfrac{\pi}{6} + k\pi \ (k \in \mathbb{Z})$}
\loigiai{
	Ta có $\tan x = \dfrac{\sqrt{3}}{3} \Rightarrow x = \dfrac{\pi}{6} + k\pi, \ k \in \mathbb{Z}$.  
	\textbf{Đáp án đúng: D}
}
\end{ex}
\begin{ex}
Trong các dãy số sau, dãy số nào là một cấp số nhân?

\choice
{$2; 4; 6; 8; 16; 32; \ldots$}
{$1; 2; 3; 4; 5; 6; \ldots$}
{$-2; -3; -4; -5; -6; -7; \ldots$}
{$1; 2; 4; 8; 16; 32; \ldots$}
\end{ex}
\begin{ex}
Công thức nghiệm của phương trình $\cos x = \cos \alpha$ là (với $k \in \mathbb{Z}$):

\choice
{$\hoac{&x = \alpha + k2\pi \\ &x = \pi - \alpha + k2\pi}$}
{$\hoac{&x = \alpha + k2\pi \\ &x = -\alpha + k2\pi}$}
{$\hoac{&x = \alpha^\circ + k360^\circ \\ &x = 180^\circ - \alpha + k360^\circ}$}
{$\hoac{&x = \alpha^\circ + k360^\circ \\ &x = -\alpha + k360^\circ}$}
\end{ex}
\begin{ex} 
Tìm $\lim \dfrac{\left(4n+2\right)\left(3n+2 \right)^4  }{ \left(3n-7 \right)^3\left(-9n^2+2n+7 \right)  }$ 

\choice 
{$- \dfrac{1}{3}$}
{\True  $- \dfrac{4}{3}$}
{$0$}
{$+\infty$ }
\loigiai{ 
	$\lim \dfrac{\left(4n+2\right)\left(3n+2 \right)^4  }{ \left(3n-7 \right)^3\left(-9n^2+2n+7 \right)  }
	= \lim \dfrac{\left[ n\left(4+\dfrac{2}{n}\right)\right]\left[n\left(3+\dfrac{2}{n}\right) \right]^4  }
	{ \left[n\left(3+\dfrac{-7}{n} \right) \right]^3 \left[n^2\left(  -9+\dfrac{2}{n}+\dfrac{7}{n^2} \right) \right]  }$\\
	$= \lim \dfrac{n^5\left[ \left(4+\dfrac{2}{n}\right)\right]\left[\left(3+\dfrac{2}{n}\right) \right]^4  }
	{ n^5\left[\left(3+\dfrac{-7}{n} \right) \right]^3 \left[\left(  -9+\dfrac{2}{n}+\dfrac{7}{n^2} \right) \right]  }
	= \lim \dfrac{\left[ \left(4+\dfrac{2}{n}\right)\right]\left[\left(3+\dfrac{2}{n}\right) \right]^4  }
	{ \left[\left(3+\dfrac{-7}{n} \right) \right]^3 \left[\left(  -9+\dfrac{2}{n}+\dfrac{7}{n^2} \right) \right]  }$
	$=\dfrac{\left(4+0\right)\left(3+0\right)^4  }{ \left(3+0 \right)^3\left(  -9+0+0 \right)  }$\\
	$=- \dfrac{4}{3}$
}
\end{ex}
\textbf{PHẦN II. Câu trắc nghiệm đúng sai}\setcounter{ex}{0}}
\begin{ex}
Cho phương trình lượng giác $2\sin x - 1 = 0 \, (*)$.

\choiceTF
{Phương trình $(*)$ tương đương $\sin x = \dfrac{1}{2}$}
{Phương trình $(*)$ tương đương $\sin x = \sin\left(\dfrac{\pi}{3}\right)$}
{Phương trình có nghiệm là $x = \dfrac{\pi}{6} + k2\pi$ hoặc $x = \dfrac{5\pi}{6} + k2\pi, \ (k \in \mathbb{Z})$}
{Trong đoạn $[0;4\pi]$ phương trình có $8$ nghiệm}
\end{ex}
\begin{ex}
Cho cấp số cộng $(u_n)$ có số hạng đầu $u_1 = \dfrac{3}{2}$, công sai $d = \dfrac{1}{2}$.  
Khi đó các mệnh đề dưới đây đúng hay sai:

\choiceTF
{Công thức số hạng tổng quát của cấp số cộng đã cho là $u_n = 1 + \dfrac{n}{2}$}
{Tổng $12$ số hạng của cấp số cộng đã cho là $144$}
{Số $\dfrac{15}{4}$ là một số hạng của cấp số cộng đã cho}
{Tổng $u_4 + u_6 = 6$}
\end{ex}
\begin{ex}
Cho biết $\sin\alpha = \dfrac{4}{5}$ với $\dfrac{\pi}{2} < \alpha < \pi$. Khi đó:
\choiceTF
{$\cos\alpha < 0$}
{$\cos\alpha = -\dfrac{3}{25}$}
{$\tan\alpha = -\dfrac{4}{5}$}
{$\cos\left(\dfrac{\pi}{3} - \alpha\right) = \dfrac{\sqrt{2}}{10}$}
\end{ex}
\begin{ex}
Cho cấp số nhân $(u_n)$ thỏa mãn $
\heva{
	&	u_3 = 2\\
	&	u_6 = 54.
}$.	Các mệnh đề sau đúng hay sai?

\choiceTF
{Tổng của $7$ số hạng đầu của cấp số nhân đã cho bằng $243$}
{$q=4$}
{Số $1458$ là số hạng thứ $9$ của cấp số nhân}
{Số hạng $u_1 = \dfrac{2}{3}$}
\end{ex}

 \textbf{PHẦN III. Câu trắc nghiệm trả lời ngắn}\setcounter{ex}{0}}
\begin{ex}
Cho $\sin\alpha = \dfrac{2\sqrt{2}}{3}$, $\cos\alpha = -\dfrac{1}{3}$.  
Khi đó, giá trị  $
\dfrac{\tan\alpha - \cot\alpha}{\tan\alpha + \cot\alpha}$
bằng bao nhiêu? (kết quả làm tròn đến hàng phần trăm)
\end{ex}
\begin{ex}
Giả sử khi một cơn sóng biển đi qua một cái cọc ở ngoài khơi, chiều cao của nước được mô hình hoá bởi hàm số 
$ h(t) = 90\cos\left(\dfrac{\pi}{3}t\right), $
trong đó $h(t)$ là độ cao tính bằng centimét trên mực nước biển trung bình tại thời điểm $t$ giây $(t \ge 0)$.  
Tìm tất cả các thời điểm trong khoảng 9 giây đầu tiên để chiều cao của sóng đạt $45$cm.  
%	\shortans[oly]{}
\end{ex}
\begin{ex}  Cho góc lượng giác $\alpha$ thỏa mãn $\sin\alpha = \dfrac{1}{3}$, $\dfrac{\pi}{2} < \alpha < \pi$.  
Tính $\sin 2\alpha$. (làm tròn đến hàng phần mười)

\loigiai{
	Vì $\dfrac{\pi}{2} < \alpha < \pi$ nên $\cos\alpha < 0$.  
	Ta có:
	\[
	\cos\alpha = -\sqrt{1 - \sin^2\alpha} = -\sqrt{1 - \dfrac{1}{9}} = -\dfrac{2\sqrt{2}}{3}.
	\]
	Khi đó:
	\[
	\sin 2\alpha = 2\sin\alpha\cos\alpha = 2 \times \dfrac{1}{3} \times \left(-\dfrac{2\sqrt{2}}{3}\right) = -\dfrac{4\sqrt{2}}{9}.
	\]
	\textbf{Kết luận:} $\sin 2\alpha = -\dfrac{4\sqrt{2}}{9}$.
}
\end{ex}
\begin{ex}
Cho cấp số cộng $(u_n)$ có
$
\heva{
	&u_1 - 2u_5 + u_6 = -15 \\
	&u_3 + u_7 = 46.
}
$.		Tìm số hạng thứ $100$ của cấp số cộng trên.

\loigiai{
	Công thức tổng quát của cấp số cộng: $u_n = u_1 + (n-1)d$.
	
	Thay vào hệ đã cho:
	\[
	\heva{
		&u_1 - 2(u_1 + 4d) + (u_1 + 5d) = -15,\\
		&(u_1 + 2d) + (u_1 + 6d) = 46.
	}
	\]
	
	Rút gọn:
	\[
	\heva{
		&u_1 - 2u_1 - 8d + u_1 + 5d = -15,\\
		&2u_1 + 8d = 46.
	}
	\Rightarrow
	\heva{
		&-3d = -15,\\
		&2u_1 + 8d = 46.
	}
	\Rightarrow
	\heva{
		&d = 5,\\
		&u_1 = 3.
	}
	\]
	
	Khi đó:
	\[
	u_{100} = u_1 + 99d = 3 + 99 \times 5 = 498.
	\]
	
	\textbf{Kết luận:} Số hạng thứ $100$ của cấp số cộng là $u_{100} = 498$.
}
\end{ex}
\begin{ex}
Sản lượng điện tháng $01/2023$ của Thủy điện Hòa Bình là xấp xỉ $1$ tỷ kWh.  
Do thời tiết $6$ tháng đầu năm $2023$ khô hạn, lưu lượng nước về hồ chứa thấp nên sản lượng điện từ tháng $01/2023$ đến tháng $06/2023$, mỗi tháng giảm $15\%$ so với tháng trước đó.  
Khi đó sản lượng điện (tỷ kWh) tháng $06/2023$ của Thủy điện Hòa Bình xấp xỉ bằng bao nhiêu? (làm tròn đến hàng phần trăm)



\loigiai{
	Đây là dãy **cấp số nhân** với:
	\[
	u_1 = 1, \quad q = 1 - 0{,}15 = 0{,}85.
	\]
	Sản lượng tháng $06/2023$ là:
	\[
	u_6 = u_1 \cdot q^{6-1} = 1 \times 0{,}85^5 \approx 0{,}44.
	\]
	\textbf{Kết luận:} Sản lượng điện tháng $06/2023$ xấp xỉ $0{,}44$ tỷ kWh.
}
\end{ex}
\begin{ex}
Cho $\sin x = \dfrac{1}{3}$. Giá trị của $\cos2x = \dfrac{m}{n}$ với $\dfrac{m}{n}$ là phân số tối giản.
Tính giá trị biểu thức $P=8m-4n$
%		\shortans[oly]{Ta có: $\cos2x = 1-2\sin^2x = 1-2\times \left(\dfrac{1}{3}\right)^2 = \dfrac{7}{9}$\\
	%			Vậy $m=7;n=9$. Vậy $P=8m-4n = 20$ }
%		\loigiai{$20$}
\end{ex}
